\subsection{Residual Output}
\label{SEC:Files:Residual_Output}

%%%%%%%%%%%%%%%%%%%%%%%%%%%%%%%%%%%%%%%%%%%%%%%%%%%%%%%%%%%%%%%%%%%%%%%%%%%%%%%%

During a run job, if the current iteration is a multiple of
\inp[nmlOutput::ResidualInterval]{Residual\_Interval},
convergence information for that iteration is written to the
\file{residual.dat} file.
%
This file is written in Tecplot ASCII format, with columns of
data and intuitive headers.
%
It contains detailed convergence information for each of the
transport equations being solved.
%
The residual measures the change in a solution variable over
the given time step and is computed as the L2 sum over all
of the solution points.

%%%%%%%%%%%%%%%%%%%%%%%%%%%%%%%%%%%%%%%%%%%%%%%%%%%%%%%%%%%%%%%%%%%%%%%%%%%%%%%%

\SeeAlso{
\inp[nmlOutput::ResidualInterval]{Residual\_Interval}
}

%%%%%%%%%%%%%%%%%%%%%%%%%%%%%%%%%%%%%%%%%%%%%%%%%%%%%%%%%%%%%%%%%%%%%%%%%%%%%%%%

